(1) 取扰动 \(e^{i(kx-\omega t)}\)，\(-i\omega n_{i1}+ikn_0v_1=0\)，\(-i\omega m_iv_1=-ike\Phi\)，\(n_{e1}=n_0e\Phi/(k_BT_e)\)。Poisson方程 \(\varepsilon_0k^2\Phi=e(n_{i1}-n_{e1})\)，符号由 \(\rho=e(n_i-n_e)\) 决定。
(2) 消去 \(v_1,n_{i1},n_{e1}\) 得 \(\omega^2=k^2c_s^2/(1+k^2\lambda_D^2)\)，其中 \(c_s^2=k_BT_e/m_i\)，\(\lambda_D^2=\varepsilon_0k_BT_e/(n_0e^2)\)。这是假定冷离子、等温电子及静电近似的离子声支。
(3) \(v_{\rm ph}=c_s(1+k^2\lambda_D^2)^{-1/2}\)，\(v_g=c_s(1+k^2\lambda_D^2)^{-3/2}\)。\(k\lambda_D\ll1\) 时均趋 \(c_s\)；\(k\lambda_D\gg1\) 时 \(\omega\to c_s/\lambda_D=\sqrt{n_0e^2/(\varepsilon_0m_i)}\)，为离子等离子体频率。
(4) 复振幅约定下 \(\overline K_i=m_in_0|v_1|^2/4=n_0e^2k^2|\Phi|^2/(4m_i\omega^2)\)。场能 \(\overline U_E=\varepsilon_0k^2|\Phi|^2/4\)，电子压缩自由能 \(\overline F_e=n_0e^2|\Phi|^2/(4k_BT_e)\)。利用色散关系得到 \(\overline K_i=\overline U_E+\overline F_e\)，正定且量纲均为能量密度。
